\documentclass{article}
\usepackage[letterpaper,top=2cm,bottom=2cm,left=3cm,right=3cm,marginparwidth=1.75cm]{geometry}

\usepackage{graphicx} % Required for inserting images
\usepackage{booktabs}
\usepackage{caption}
\usepackage{subcaption}
\usepackage{authblk}
\usepackage{amsmath}
\usepackage{amssymb}
\usepackage[dvipsnames]{xcolor}
\usepackage[english]{babel}
\usepackage{pgfgantt}
\usepackage{biblatex}
\bibliography{latex/bibliography}
\newcommand{\mycitethis}{\tkDE{[?]}}
\usepackage[colorlinks=true, allcolors=blue]{hyperref}

\usepackage{verbatim}

\newcommand{\detailtexcount}[1]{%
  \immediate\write18{texcount -merge -sum -q main.tex output.bbl > #1.wcdetail }%
  \verbatiminput{#1.wcdetail}%
}

% Change font
\usepackage{helvet}
\renewcommand{\familydefault}{\sfdefault}

\usepackage{xcolor}
\newcommand{\tkDE}[1]{\textcolor{magenta}{#1}} 
\newcommand{\citethis}{\tkDE{[?]}}

\newcommand{\diff}{\mathrm{d}}  

\title{Protocol: Ethnicity-based System Dynamics Model of Measles Vaccinations in Birmingham and Solihull}
\author{David Ellis}

\date{}

\begin{document}
\maketitle

\section{Background}

Measles is a highly infectious viral disease that can have serious complications and fatality \cite{WHO_measles_topic}. 
In high-income countries, complications occur in around 10-20\% of cases and the case fatality ratio is around 0.05\% in  \cite{wolfson2009estimates, NICE2024MeaslesComplications}. Clinically vulnerable groups, such as immunosuppressed individuals, infants under one year, and pregnant people, are more likely to develop complications \cite{NHSEngland2024MeaslesGuidance}. 

The first measles vaccine was licensed in the US in 1963 \cite{MOSS20172490}. The measles vaccine was first introduced in the UK in 1968, later replaced by the combined measles, mumps, and rubella (MMR) vaccine in 1988 \cite{ukhsa_vaccination_timeline_1796_present}. As of 1 January 2026, varicella has been added to make the MMRV vaccine. In the UK all children are offered a first dose of MMRV when they turn 12 months and a second dose at 18 months \cite{UKHSA2023MeaslesRubellaElimination}. The World Health Organisation (WHO) estimates that between 2000 and 2024, the measles vaccine has prevented  59 million deaths \cite{mondiale2025weekly}. 

The UK achieved measles elimination in 2016 \cite{UKHSA2023MeaslesRubellaElimination}. However, measles outbreaks are again becoming increasing common in the UK as well as globally due to falling vaccine coverage rates \cite{minta2023progress}. The top half of Figure.~\ref{fig:measles-background} shows the that there has been a series of measles outbreaks in England since January 2012. 

\begin{figure}[ht!]
    \centering
    \includegraphics[width=0.7\linewidth]{figures/background/measles-background.pdf}
    \caption{ (Top) Monthly lab-confirmed cases of measles in England \cite{UKHSA2025MMRLabConfirmedCasesQ3}. (Bottom) Yearly MMR coverage for one dose at two years old \cite{FTips}.}
    \label{fig:measles-background}
\end{figure}


The bottom half of Figure.~\ref{fig:measles-background} shows that these outbreaks have occurred in a context of falling vaccination coverage nationally from a peak of 92.7\% in 2013 down to 88.9\% in 2024.  However, in Birmingham and Solihull, the  coverage started lower and has fallen further than the England average from a peak of 89.7\% in 2014 down to 81.3\% in 2024 \cite{FTips}. This has contributed to the region suffering large measles outbreaks \cite{jary2025sociodemographic}. 

Part of the reason for the falling MMR coverage is the growing scepticism in vaccines \citethis. While this is a trend for most vaccines, there has been particular public concern about the MMR vaccine due to the infamous, since retracted, claim by one research group that the MMR vaccine might cause autism \cite{novilla2023parents, wakefield1998retracted, godlee2011wakefield}. However, there many factors leading to vaccine hesitancy that impact different groups differently including misinformation, health literacy, and lack of trust in the healthcare system \cite{cagnotta2025parental}. Some groups, such as Gypsy, Roma, and Traveller populations, also face additional barriers to accessing the vaccine \cite{wilder2020resurgence}. As a result, significant ethnic and socioeconomic inequality has been observed in MMR coverage \cite{Flatte079550, zhang2023ethnic}. 

During the outbreak in Birmingham from 13 October 2023 the measles infection rate was 3.4 times higher for those living in the most deprived quintile compared to those living in the third most deprived quintile. No infections occurred in the least deprived quintile. Additionally, the infection rate for Black African residents was eight times greater than the white British population \cite{Tabacchi02072016}. Figure.~\ref{fig:BSol-mmr-eth} shows that, except for those of Indian heritage, children off all ethnic groups have lower MMR coverage at one dose than the White British cohort. The Black African and Black Caribbean have particularly low coverage at 74.2\% and 64.0\% respectively compared to 92.9\% for Whit British \cite{NHSEngland2026FDP}. 

\begin{figure}[ht!]
    \centering
    \includegraphics[width=0.91\linewidth]{figures/background/mmr-background.pdf}
    \caption{Percentage of children aged less than six years old living in Birmingham and Solihull who have received at least one dose of MMR or MMRV as of 04/02/2026 \cite{NHSEngland2026FDP}. The error bars indicate the 95\% CI calculated using the Wilson method.}
    \label{fig:BSol-mmr-eth}
\end{figure}

Among the Core Cities of England and Wales\footnote{England and Wales Core Cities include Birmingham, Bristol, Cardiff, Leeds, Liverpool, Manchester, Newcastle, Nottingham, and Sheffield.}, Birmingham has the highest level of ethnic diversity. Data from the 2021 Census show that just under half (48.5\%) of Birmingham and Solihull residents identified as White British, meaning that a majority (51.5\%) reported a different ethnic background. This is almost double the England average of 26.5\%. The largest ethnic minority groups were Pakistani (14.9\%), Indian (5.7\%), Black African (5.0\%), and Other White (3.8\%) \cite{ONS2021census}.


All six regions of the World Health Organization (WHO) have declared measles elimination goals \cite{WHO_measles_topic}. The UK Measles and Rubella elimination strategy 2019 set the MMR vaccination target, based on WHO recommendation, to 95\% of two doses by 5 years of age in order to achieve and maintain measles elimination \cite{public2019uk}. The Public Health England (PHE) immunisation inequalities local action plan template outlines that factors responsible for inequitable vaccine uptake are context specific \cite{PHE2021ImmunisationInequalities}. Therefore, in order to reverse the falling measles vaccine coverage and growing health inequalities, it is therefore imperative that we understand the range of factors influencing parent's decision to vaccinate their children and how these factors are influenced by ethnicity. 

This project aims to build a decision-support system in collaboration with a range of local health protection experts working in Birmingham and Solihull. An interactive workshop will be hosted to identify behavioural and epidemiological factors that facilitate and block flows in the system and to identify potential interventions to improve these flows. This model, including selected interventions, will be built using a system dynamics approach.

\section{Literature Review}

This project requires a three pronged literature review:
\begin{itemize}
    \item Modelling measles epidemics including vaccination (in UK/Europe?)
    \item Factors influencing parents decisions to vaccinate their children
    \item Interventions used to increase child vaccinations
\end{itemize}

We will adopt a semi-structured approach finding literature. We will use Medline and Proquest and supplement with additional papers found through other means. Local intelligence will also be identified through consultation with health protection leads at BCC and Solihull Metropolitan Council. English language papers published in the last 10 years (2016 or more recent).

\subsection{Modelling Measles Epidemics}
(vaccine OR vaccination) AND measles AND simulation

\tkDE{Focus on why not other methods.}

We see that, to date, there has not been a attempt to simulate a measles vaccination system that considers the interaction between epidemiological and behavioural factors including the variation in these factors across ethnicity.

\tkDE{Prisma flow diagram (optional extra)}

\subsection{Factors influencing vaccine decisions}
Specifically related to parental decisions to vaccinate their children. 

(vaccine or vaccination or immunisation) AND (barrier OR hesitancy) AND (child OR parent) AND (measles OR MMR)

PubMed - 16 results (specifically filtered for systematic reviews and meta analyses)
ProQuest - 35 results

One study excluded since it related specifically to conflict setting.
% https://link.springer.com/article/10.1007/s10900-014-9956-1
%  PHE immunisation health equity audit includes an evidence-based model which describes the factors influencing inequality in vaccination uptake in high income settings
Papers included are presented in Table.~\ref{tab:vaccine-hesitancy-papers}.


%\tkDE{Prisma flow diagram (optional extra)}

\subsection{MMR uptake interventions}
What are the barriers to vaccinations?

intervention AND vaccination AND MMR

% Evaluating the impact of national and regional measles catch-up activity on MMR vaccine coverage in England, 2023 to 2024 https://www.gov.uk/government/publications/evaluation-of-vaccine-uptake-during-the-2023-to-2024-mmr-catch-up-campaigns-in-england/evaluating-the-impact-of-national-and-regional-measles-catch-up-activity-on-mmr-vaccine-coverage-in-england-2023-to-2024

\section{Research Objectives}

This project will aim to answer the following research question: How can a system dynamics model be used to facilitate the development of interventions to increase measles vaccination rates and reduce negative outcomes, particularly amongst lower-uptake ethnic groups?

The research aims are as follows:
\begin{itemize}
    \item Conduct a participatory modelling workshop with expert stakeholders to design a theoretical model for measles vaccinations and outcomes in Birmingham and Solihull.
    \item Identify and prioritise key interventions to be modelled
    \item Develop a functional model and calibrate using historical vaccination, notification, and hospitalisation data
    \item Integrate most highly prioritised interventions to simulate the expected impact on the modelled behavioural science factors, vaccination rates, and health outcomes
\end{itemize}

Long term, the aim is to develop a dashboard for health protection teams working in the system to be able to dynamically plan and test interventions using variable model assumptions. However, this is beyond the scope of the thesis. 

\section{Methods}

\subsection{Introduction to System Dynamics}
Following approach from \cite{van2020explaining} using Stella Architect. Stocks and flows. \tkDE{Include example diagram?}

\begin{itemize}
    \item Clear explanation of what SD and Stella Architect are. 
    \item Assumptions
    \item Limitations
\end{itemize}

\subsection{Stakeholder Engagement}

Following approach from Michail et. al \cite{michail2025evaluation}. 

Currently (\tkDE{26/02/2026}), \tkDE{eight}, stakeholders have agreed to participate in the modelling workshop. Stakeholder organisations include BCC (5), Solihull Metropolitan Council (1), and BSol NHS ICB (2). To help promote full participation, stakeholders will be separated into \tkDE{two} groups. The make-up of each group will be chosen to create the best spread of background and expertise. For example, two of the participants from the BSol ICB have prior experience in system dynamics and will therefore be split up to maximise the conceptual support given to each group. 

Participants will be sent an evidence pack outlining the rational behind using a system dynamics approach rather than a more traditional outbreak simulation technique as well as published and local evidence on the systemic, sociological, and psychological barriers that prevent parents from vaccinating their children. 

Input from the participants will be collected on paper and sticky notes. At the end of the session, all materials will be photographed and later transcribed. For transparency, the photos and transcriptions will be made publicly available via the GitHub repository. The session will be facilitated with the help of the BCC Governance Team. 

\subsubsection{Part 1 - Introduction}
\begin{itemize}
    \item High-level overview of the evidence pack
    \item Practical model building task to help participants understand the basics behind system dynamics using an an normal-life example that everyone can equally relate to and engage with (probably the banking example)
    \end{itemize}

\subsubsection{Part 2 - Model building}
\begin{itemize}
    \item Stakeholders will be presented with a simplified measles "model" showing the main stocks (states) in the system (\tkDE{Outline and probably include a diagram}). For each of the key flows between stocks, participants will be asked to use sticky notes to give factors that facilitate and block the flow. Once the factors are listed, participants will be given a limited number of stickers each to vote on the factors for each flow that they think are dominant. While all factors are important, this will help to focus the model on the factors with the biggest influence on the model's behaviour.
\end{itemize}

\subsubsection{Part 3 - Interventions}
\begin{itemize}
    \item Participants will be asked to identify which flows they think are problematic (Assuming given model is complex enough to justify this step).
    \item Participants will then be asked to identify policies and interventions that could change this behaviour and, for each to identify clearly who will be doing what, where in the system this will occur, and what the expected direction of change is.
    \item This information will also be collected on sticky notes. 
    \item Each group will then choose two or three interventions to share with the group. Once all the interventions are collected, Menti will then be used to allow every participant to rank each intervention by preference. This will help to focus the project on the interventions that stakeholders think are the most worth focusing on based on their feasibility and impracticability. 
\end{itemize}

\subsection{Model Development}
\begin{itemize}
    \item Basic model development
    \item Take back to stakeholders for validation
    \item Calibration
    \item Intervention implementation
    \item Feed back to stakeholders as required
\end{itemize}

\subsection{Data Sources}

\begin{itemize}
    \item Population data (GP registrations) – BSol ICB
    \item Child MMR vaccinations – Federated data platform
    \item Measles lab-confirmed notifications – UKHSA (FOI)
    \item Hospitalisation data – BSol ICB
\end{itemize}

\subsection{Data Analysis}
GP registration and SUS data will be extracted using SQL. All data analysis and visualisation will be completed in R.

\subsection{Sensitivity Analysis}

\section{Feasibility}

\begin{itemize}
    \item Data completeness
    \item Ethical approval is not required since the data used is pseudonymised and routinely collected and the purpose of this project is for service improvement so this will not impact the project feasibility. 
    \item Talk about training and experience with R, SQL, and Stella?
\end{itemize}


\section{Timeline}

% include thesis
\hspace*{-2cm}
\begin{ganttchart}[%Specs
     x unit = .3cm,  %<---------------------- New x unit 
     y unit title=0.5cm,
     y unit chart=0.7cm,
     vgrid, hgrid,
     title height=1,
%     title/.style={fill=none},
     title label font=\bfseries\footnotesize,
     bar/.style={fill=blue},
     bar height=0.7,
%   progress label text={},
     group right shift=0,
     group top shift=0.7,
     group height=.3,
     group peaks width={0.2},
     milestone/.style={fill=magenta},
     inline]{1}{36}
    %labels

    \gantttitle[]{Jan}{4}
    \gantttitle[]{Feb}{4}
    \gantttitle[]{Mar}{4}
    \gantttitle[]{Apr}{4}
    \gantttitle[]{May}{4}
    \gantttitle[]{Jun}{4}
    \gantttitle[]{Jul}{4}
    \gantttitle[]{Aug}{4}
    \gantttitle[]{Sep}{4}
    
    % title
   % \gantttitle{Q1}{1}                  
    \\
    

    % Setting group if any
     %\ganttgroup[inline=false]{Phase 1: Orientation}{1}{4}\\
     \ganttgroup[inline=false]{Phase 1: Planning}{1}{10}\\
    \ganttbar[inline=false]{Project Scoping}{1}{4}\\  

    \ganttbar[inline=false]{Write Protocol Presentation}{4}{7}\\  
    \ganttmilestone[inline=false]{Present Protocol}{7}\\
    \ganttbar[inline=false]{Write Protocol}{7}{10}\\  
    \ganttmilestone[inline=false]{Submit Protocol}{10}\\
    
  
    \ganttgroup[inline=false]{Phase 2: Participatory Model Designing}{4}{11}\\
    \ganttbar[inline=false]{Stakeholder Mapping}{4}{7}\\  
    \ganttmilestone[inline=false]{Send workshop invites}{7}\\
    
    \ganttbar[inline=false]{Develop Evidence Pack}{7}{10}\\
    \ganttmilestone[inline=false]{Modelling Workshop}{11}\\
   
    \ganttgroup[inline=false]{Phase 3: Model Development}{11}{21}\\
    \ganttbar[inline=false]{Preliminary modelling building}{11}{14}\\
    \ganttbar[inline=false]{Calibrate model}{15}{16}\\
    \ganttmilestone[inline=false]{Present model to stakeholders}{16}\\   
    \ganttbar[inline=false]{Implement stakeholder feedback}{16}{17}\\
    \ganttbar[inline=false]{Model interventions}{17}{20}\\
    
    \ganttgroup[inline=false]{Phase 3: Data Analysis}{21}{24}\\
    \ganttbar[inline=false]{Create visualisations in R}{21}{23}\\
    \ganttbar[inline=false]{Sensitivity Analysis}{22}{24}\\
    \ganttmilestone[inline=false]{Present preliminary results}{25}\\   
    
    \ganttgroup[inline=false]{Phase 3: Thesis writing}{25}{33}\\
    \ganttbar[inline=false]{Write first draft}{25}{29}\\
    \ganttmilestone[inline=false]{Submit draft}{29}\\
    \ganttbar[inline=false]{Implement feedback}{29}{33}\\
    \ganttbar[inline=false]{Design Poster}{30}{33}\\
    \ganttmilestone[inline=false]{Present Poster}{33}\\   
    \ganttmilestone[inline=false]{Submit final thesis}{33}
    
\end{ganttchart}

%\bibliographystyle{unsrt}
%\bibliography{bibliography}

%TC:ignore
\printbibliography

\newpage
\appendix
\section{Literature Search}

\begin{table}[ht!]
    \centering
    \begin{tabular}{c|c}
         &  \\
         & 
    \end{tabular}
    \caption{Papers on epidemic modelling.}
    \label{tab:placeholder}
\end{table}


\begin{table}[ht!]
\begin{tabular}{p{0.6\linewidth} | p{0.15\linewidth} | p{0.06\linewidth} | p{0.1\linewidth}}
\toprule
\textbf{Title} & \textbf{Authors} & \textbf{Year} & \textbf{Citation} \\ \midrule

\citetitle{novilla2023parents} & \citeauthor*{novilla2023parents} & \citeyear{novilla2023parents}  & \cite{novilla2023parents} \\ \midrule

\citetitle{essa2024understanding} & \citeauthor*{essa2024understanding} & \citeyear{essa2024understanding}  & \cite{essa2024understanding} \\ \midrule

\citetitle{maamor2024prevalence} & \citeauthor*{maamor2024prevalence} & \citeyear{maamor2024prevalence}  & \cite{maamor2024prevalence} \\ \midrule

\citetitle{wilder2020resurgence} & \citeauthor*{wilder2020resurgence} & \citeyear{wilder2020resurgence}  & \cite{wilder2020resurgence} \\ \midrule

\citetitle{crawshaw2022defining} & \citeauthor*{crawshaw2022defining} & \citeyear{crawshaw2022defining}  & \cite{crawshaw2022defining} \\ \midrule

\citetitle{ames2017parents} & \citeauthor*{ames2017parents} & \citeyear{ames2017parents}  & \cite{ames2017parents} \\ \midrule

\citetitle{halliday2025vaccine} & \citeauthor*{halliday2025vaccine} & \citeyear{halliday2025vaccine}  & \cite{halliday2025vaccine} \\ \midrule

\citetitle{bocquier2017socioeconomic} & \citeauthor*{bocquier2017socioeconomic} & \citeyear{bocquier2017socioeconomic}  & \cite{bocquier2017socioeconomic} \\ \midrule

\citetitle{baiocchi2025parental} & \citeauthor*{baiocchi2025parental} & \citeyear{baiocchi2025parental}  & \cite{baiocchi2025parental} \\ \midrule

\citetitle{gschwend2026shadow} & \citeauthor*{gschwend2026shadow} & \citeyear{gschwend2026shadow}  & \cite{gschwend2026shadow} \\ \midrule

\citetitle{cagnotta2025parental} & \citeauthor*{cagnotta2025parental} & \citeyear{cagnotta2025parental}  & \cite{cagnotta2025parental} \\ \midrule

\citetitle{tankwanchi2021vaccine} & \citeauthor*{tankwanchi2021vaccine} & \citeyear{tankwanchi2021vaccine}  & \cite{tankwanchi2021vaccine} \\ \midrule

\citetitle{torracinta2021mmr} & \citeauthor*{torracinta2021mmr} & \citeyear{torracinta2021mmr}  & \cite{torracinta2021mmr} \\ \midrule

\citetitle{Tabacchi02072016} & \citeauthor*{Tabacchi02072016} & \citeyear{Tabacchi02072016}  & \cite{Tabacchi02072016} \\ 

\bottomrule
\end{tabular}
    \caption{Papers on vaccine hesitancy.}
    \label{tab:vaccine-hesitancy-papers}
\end{table}

\begin{table}[ht!]
    \centering
    \begin{tabular}{c|c}
         &  \\
         & 
    \end{tabular}
    \caption{Papers on child-vaccine interventions.}
    \label{tab:placeholder}
\end{table}

\end{document}
